%% appD-terminology \dash{} Słownik terminów
%
% Wiersze sekcji D.3 i D.4 to \plterm{}{}{}, co jest celowo greppowalne:
% `check_structure.py --terms` wymaga, by każde polskie tłumaczenie wymienione
% tutaj występowało w prozie programs/pl. Słownik to na każdej linii twierdzenie
% o treści książki, a usterka, która naprawdę boli, to wiersz nazywający słowo,
% którego książka nie używa \dash{} a proponowana tabela w notes/03 miała takie
% dwa.
%
% Argumenty 1 i 2 są identyczne w obu wydaniach: hasłem jest termin angielski,
% bo to jego czytelnik będzie szukał. Tłumaczona jest wyłącznie notatka.
%
% Tabela w sekcji D.1 NIE jest bramkowana i jest pisana ręcznie celowo. Obie
% jej kolumny to dosłowny tekst, bo wpisanie tam makra kontraktu wypisałoby
% formę bieżącego wydania w OBU kolumnach i nie powiedziałoby niczego -- a poza
% tym C10 zakazuje \tan, \cot, \arctan, \gcd i \operatorname{lcm} wprost i
% słusznie, bo proza nigdy nie może zaszywać formy, którą drugie wydanie składa
% inaczej. Parytet C3 pilnuje, by oba pliki językowe definiowały ten sam ZBIÓR
% makr; nic nie sprawdza, czy ta tabela opisuje je poprawnie. Czytaj ją więc
% razem z lang/en.tex i lang/pl.tex, ilekroć któryś się zmieni.

\chapter{Słownik terminów}
\label{app:D}

\noindent
Dwa języki, jedna książka i jedna strona, na której spisano wybory. Ten dodatek
jest dla czytelnika obu wydań: polskiego, który musi wiedzieć, za jakie
angielskie słowo stoi dany polski termin, i angielskiego, który może czytać
polski artykuł albo siedzieć na polskim przeglądzie projektu.

Mówi też, gdzie wydanie polskie \emph{nie} ustaliło jednego słowa. To trudniejsza
połowa i jest tutaj, a nie sprzątnięta pod dywan, bo czytelnik, który napotyka
dwa słowa na jedną rzecz, musi usłyszeć, że to jedna rzecz.

\section{Symbole, które oba wydania składają inaczej}
\label{sec:D-notation}
\index{notation!contract}
\index{terminology!notation contract}

Reguła to \textbf{test klawiatury}: jeśli czytelnik wpisze token do komputera,
oba wydania używają zapisu kodu; jeśli będzie go tylko czytał albo pisał
ręcznie, wydanie polskie używa formy polskiej. Dlatego wydanie polskie pisze
\emph{tg} na tangens i $\tanh$ \dash{} a nie \emph{tgh} \dash{} na tangens
hiperboliczny, i ta pozorna niekonsekwencja jest regułą działającą poprawnie,
a nie przeoczeniem.

Każdy wiersz poniżej to jedno makro w \code{lang/en.tex} i \code{lang/pl.tex},
więc \emph{źródło} obu wydań jest identyczne i różni się tylko wyjście. Makro
zdefiniowane w jednym pliku językowym i nie w drugim to niezdefiniowana
sekwencja sterująca w dokładnie jednej z dwóch kompilacji, i dlatego ta para
jest bramkowana.

Obie kolumny są złożone jako dosłowny tekst i jest to celowe. Wpisanie tu makra
kontraktu wypisałoby formę bieżącego wydania w obu kolumnach i nie powiedziałoby
niczego.

\begin{center}
\begin{tabularx}{\linewidth}{@{}llX@{}}
\toprule
\textbf{Wydanie angielskie} & \textbf{Wydanie polskie} & \textbf{makro i co to jest} \\
\midrule
\emph{tan}    & \emph{tg}     & \code{\textbackslash tg}, tangens \\
\emph{cot}    & \emph{ctg}    & \code{\textbackslash ctg}, cotangens \\
\emph{arctan} & \emph{arc tg} & \code{\textbackslash arctg}, tangens odwrotny \\
\emph{arccot} & \emph{arc ctg} & \code{\textbackslash arcctg}, cotangens odwrotny \\
\emph{gcd}    & \emph{NWD}    & \code{\textbackslash gcdop}, największy wspólny dzielnik \\
\emph{lcm}    & \emph{NWW}    & \code{\textbackslash lcmop}, najmniejsza wspólna wielokrotność \\
\emph{span}   & \emph{lin}    & \code{\textbackslash spanof}, powłoka liniowa zbioru wektorów \\
0.5           & 0,5           & \code{\textbackslash num} i \code{\textbackslash val}, znak dziesiętny \\
\bottomrule
\end{tabularx}
\end{center}

\medskip
I kilka, których celowo \emph{nie} ma na tej liście, bo oba wydania składają je
tak samo. $\Prob$, $\Ex$, $\Var$ i $\Cov$ pisze się tak, jak pisze się je po
angielsku, a przedział $\intcc{a}{b}$ zachowuje nawiasy kwadratowe w obu.
Wariancja i przedział to te dwa, które najpewniej zaskoczą polskiego
czytelnika, i oba wracają w sekcji~\ref{sec:D-unsettled}.

\section{Gdzie książka zatrzymuje się, by powiedzieć, że symbol znaczy co innego}
\label{sec:D-boxes}
\index{notation!collisions}

Ramka z oznaczeniami to książka zatrzymująca się w środku ramki, żeby
powiedzieć, że symbol albo słowo zaraz będzie znaczyło coś innego niż to, co
czytelnik ze sobą przynosi. Warto wypisać je w jednym miejscu, bo każda z nich
znaczy kolizję, w którą ktoś już wszedł.

\begin{center}
\begin{tabularx}{\linewidth}{@{}lX@{}}
\toprule
\textbf{gdzie} & \textbf{co rozstrzyga} \\
\midrule
\ref{prog:F02} & litera jest nazwą, a nie zagadką \\
\ref{prog:F03} & nigdy goły logarytm, i dlaczego zakaz ma wyjątki \\
\ref{prog:F04} & jedno angielskie słowo w dwóch rolach \\
\ref{prog:P01} & mantysa, i co to słowo ze sobą niesie \\
\ref{prog:P03} & logarytm bez podstawy, wewnątrz $O$ \\
\ref{prog:P12} & symbol Newtona, czytany na dwa sposoby \\
\ref{prog:P15} & \emph{gradient} w dwóch rolach: nachylenie i wektor \\
\ref{prog:P18} & układ licznikowy, zadeklarowany raz na całą książkę \\
\ref{prog:P23} & \emph{niezależny} w dwóch rolach \\
\ref{prog:P24} & wielkie i małe litery; wartość oczekiwana jest liczbą; oraz
                 $\Var(X)$ wobec $D^{2}(X)$ \\
\bottomrule
\end{tabularx}
\end{center}

\medskip
Ostatnia z nich znaczy najwięcej dla polskiego czytelnika i nie jest
opcjonalna: $D^{2}(X)$ i $M(X)$ to \emph{obecny} polski uzus dydaktyczny, a nie
historyczna ciekawostka, więc czytelnik, który się ich nauczył i napotyka
niewyjaśnione $\Var$, ma powód, by zwątpić w książkę, a nie w swoją szkołę.
Książka i tak pisze $\Var$ oraz $\Ex$, a powód jest ten, na którym stoi cały ten
dodatek: tak piszą biblioteki i artykuły, które czytelnik będzie otwierał.
Program~\ref{prog:P24} niesie tę odpowiedniość w obie strony, więc żaden
z zapisów nie zostaje sam.

\section{Terminy, hasłem angielskie}
\label{sec:D-terms}
\index{terminology!glossary}

Hasłem jest termin angielski, bo to jego czytelnik będzie szukał, czytał w
artykule i widział w API. Tam, gdzie polska forma jest w realnym użyciu, wydanie
polskie jej używa. Tam, gdzie żadnej nie ma, wydanie polskie zostawia słowo
angielskie i odmienia je po polsku \dash{} \emph{embeddingów},
\emph{tokenizator}, \emph{benchmarku} \dash{} bo tak mówią polscy inżynierowie,
a kalka wymyślona przez tłumacza jest gorsza niż zapożyczenie, którego wszyscy
już używają.

Trzy tabele to ta reguła, posortowana. Te wydanie polskie tłumaczy.

\begin{center}
\begin{tabularx}{\linewidth}{@{}llX@{}}
\toprule
\textbf{Angielski} & \textbf{Wydanie polskie} & \textbf{notatka} \\
\midrule
\plterm{activation}{aktywacja}{Ustalone. \emph{Funkcja aktywacji} na samą
  funkcję.}
\plterm{attention}{uwaga}{\emph{mechanizm uwagi}, \emph{warstwa uwagi},
  \emph{głowica uwagi}. Realny polski uzus; samego terminu angielskiego się
  nie używa.}
\plterm{backpropagation}{propagacja wsteczna}{Ustalone.
  Program~\ref{prog:F12} mówi, czym jest: regułą łańcuchową zastosowaną do
  złożenia warstw.}
\plterm{cross-entropy}{entropia krzyżowa}{Ustalone.}
\plterm{evaluation set}{zbiór ewaluacyjny}{Ustalone.}
\plterm{feature}{cecha}{Ustalone i używane wszędzie.}
\plterm{gradient descent}{spadek gradientu}{Ustalone. Książka pisze i
  \emph{spadek gradientu}, i \emph{spadek gradientowy}, co jest jedną nazwą
  w dwóch formach gramatycznych, a nie dwoma słowami.}
\plterm{layer}{warstwa}{Ustalone.}
\plterm{loss}{strata}{Ustalone. \emph{Funkcja straty} na funkcję.}
\plterm{optimiser}{optymalizator}{Ustalone.}
\plterm{overfitting}{przeuczenie}{Ustalone.}
\plterm{perplexity}{perpleksja}{Ustalone. Program~\ref{prog:P29} mówi, co ta
  liczba zlicza.}
\plterm{probe}{sonda}{Zwykłe polskie słowo, realnie używane w tym znaczeniu,
  więc nie kalka. Program~\ref{prog:P31}.}
\plterm{sequence}{sekwencja}{Ustalone.}
\plterm{vocabulary}{słownik}{Ustalone.}
\plterm{weight}{waga}{Ustalone.}
\bottomrule
\end{tabularx}
\end{center}

\medskip
A te zapożycza, bo żadna polska forma nie jest w realnym użyciu. Słowo
angielskie zostaje i odmienia się po polsku, bo tak mówią polscy
inżynierowie.
\begin{center}
\begin{tabularx}{\linewidth}{@{}llX@{}}
\toprule
\textbf{Angielski} & \textbf{Wydanie polskie} & \textbf{notatka} \\
\midrule
\plterm{benchmark}{benchmark}{Zapożyczone i odmieniane.}
\plterm{checkpointing}{checkpointing}{Zapożyczone. Wymiana arytmetyki na bajty,
  Program~\ref{prog:P16}.}
\plterm{harness}{harness}{Zapożyczone i odmieniane; \emph{zestaw} był próbowany
  i zderzył się ze \emph{zbiorem}.}
\plterm{hyperparameter}{hiperparametr}{Zapożyczone i odmieniane.}
\plterm{logit}{logit}{To samo słowo w polskiej statystyce.}
\plterm{prompt}{prompt}{Zapożyczone i odmieniane.}
\plterm{softmax}{softmax}{Zapożyczone; nigdzie nietłumaczone.}
\plterm{token}{token}{Zapożyczone i odmieniane. \emph{Tokenizator} na to, co je
  produkuje.}
\plterm{transformer}{transformer}{Zapożyczone. Nigdy \emph{transformator}, bo
  to urządzenie elektryczne.}
\bottomrule
\end{tabularx}
\end{center}

\medskip
Te wymagają czegoś więcej niż tłumaczenia. W każdym jedno słowo pełni dwie
role albo oba wydania ustaliły co innego, a czytelnik, który weźmie pierwsze
znaczenie, gdzieś w książce się pomyli.
\begin{center}
\begin{tabularx}{\linewidth}{@{}llX@{}}
\toprule
\textbf{Angielski} & \textbf{Wydanie polskie} & \textbf{notatka} \\
\midrule
\plterm{head}{głowica}{W uwadze. \emph{Głowa} to ta anatomiczna.}
\plterm{inference}{wnioskowanie}{Dwie role, a polszczyzna przenosi angielską
  kolizję, zamiast ją rozstrzygać: wnioskowanie statystyczne,
  Programy~\ref{prog:P27} i \ref{prog:P28}, oraz uruchamianie wytrenowanego
  modelu, jak w \emph{przy wnioskowaniu} w Programie~\ref{prog:F02}.}
\plterm{normalisation}{normalizacja}{Dwie role, po polsku tak samo jak po
  angielsku: dzielenie wektora przez jego długość, Program~\ref{prog:F09}, oraz
  normalizacja warstwy, Program~\ref{prog:P32}.}
\plterm{precision}{precyzja}{Dwie role, a kolizja jest w obu językach: miara
  klasyfikacji z Programu~\ref{prog:P23} oraz szerokość mantysy liczby
  zmiennoprzecinkowej w całym Programie~\ref{prog:P01}.}
\plterm{recall}{czułość}{Ustalone. Stoi obok \emph{precyzji} w tym samym
  raporcie klasyfikacji i warunkuje po drugiej populacji.}
\plterm{residual connection}{rezydualny}{\emph{Połączenie rezydualne},
  \emph{strumień rezydualny}. Program~\ref{prog:P32}.}
\plterm{residual of a fit}{reszta}{Inny obiekt, a polszczyzna daje mu inne
  słowo: to, co zostawia dopasowanie najmniejszych kwadratów,
  Program~\ref{prog:P08}. Angielski ma na oba jedno słowo.}
\plterm{significand}{mantysa}{Wydania rozstrzygają to inaczej i oba mają
  rację. Standard mówi \emph{significand} o przechowywanych bitach, a
  \emph{mantysę} rezerwuje dla części ułamkowej logarytmu, więc wydanie
  angielskie pisze \emph{significand}; polska praktyka mówi \emph{mantysa} na
  jedno i drugie, więc wydanie polskie też. Program~\ref{prog:P01} ma ramkę
  z oznaczeniami w obu wydaniach, w każdym z przeciwnym wnioskiem.}
\bottomrule
\end{tabularx}
\end{center}

\section{Gdzie polski uzus się nie ustalił}
\label{sec:D-unsettled}
\index{terminology!unsettled}

Cztery terminy w tej książce są w wydaniu polskim oddawane więcej niż jednym
słowem. Są wypisane, a nie po cichu rozstrzygnięte, bo czytelnik, który napotyka
drugie słowo, musi wiedzieć, że to pierwsze.

\begin{center}
\begin{tabularx}{\linewidth}{@{}llX@{}}
\toprule
\textbf{Angielski} & \textbf{Wydanie polskie} & \textbf{notatka} \\
\midrule
\plterm{batch}{partia, wsad, batch}{Jeden obiekt, a każdy z wypisanych obok
  odpowiedników ma swoje zadanie. \emph{Partia} jest podstawą: jest
  najczęstsza i jako jedyna tworzy złożenie, skąd bierze się
  \emph{mikropartia}. \emph{Wsad} zostaje przy przymiotniku \emph{wsadowy},
  którego \emph{partia} nie tworzy. \emph{Batch} to zapożyczenie i zostaje
  tam, gdzie czytelnik wpisze je w wyszukiwarkę. Ilu ich jest, celowo tu nie
  liczymy: policzono to raz, a pod spodem zdążył pojawić się czwarty.}
\plterm{embedding}{zanurzenie, embedding}{Własna reguła tłumaczeniowa tej
  książki każe zostawiać słowo angielskie tam, gdzie polska forma jest kalką,
  której nikt nie wyszuka, a wydanie polskie i tak używa głównie
  \emph{zanurzenia}. Powodem jest czasownik: \emph{zanurzać} nie ma formy
  zapożyczonej \dash{} nikt nie mówi \emph{embedduje} \dash{} więc
  rzeczownik poszedł za nim. \emph{Embedding} to termin, który polski inżynier
  wpisze w wyszukiwarkę.}
\plterm{learning rate}{współczynnik uczenia, krok uczenia}{Obie formy są w
  realnym użyciu. \emph{Krok} to \emph{step}, czyli dokładnie to, co ten
  współczynnik mnoży.}
\plterm{sample}{próba, próbka}{\emph{Próba} to termin statystyczny i to jego
  książka używa w utartych połączeniach \dash{} \emph{wariancja z próby},
  \emph{średnia z próby}, Program~\ref{prog:P26}. \emph{Próbka} pojawia się
  tam, gdzie próba jest konkretną rzeczą w ręku, jak w bootstrapie
  Programu~\ref{prog:P27}. \emph{Próbkowanie} to operacja w obu przypadkach.}
\bottomrule
\end{tabularx}
\end{center}

\medskip
\begin{note}
Najdotkliwszym kosztem nieustalonego słowa jest odsyłacz, który przemianowuje
rzecz, na którą wskazuje. Tabela tego, czego ta książka nie zmierzyła, w
Programie~\ref{prog:P32} odsyła czytelnika do Programu~\ref{prog:P08} i
Programu~\ref{prog:P11} po widmo prawdziwej macierzy zanurzeń; wydanie polskie
pisze tam \emph{macierz zanurzeń}, tak samo jak Program~\ref{prog:P11}, a
Program~\ref{prog:P08} pisze \emph{macierz embeddingów}. Tego nie da się
naprawić przy samym odsyłaczu, bo wiersz wskazuje dwa programy, które nie
zgadzają się ze sobą, i to jest argument za rozstrzygnięciem \emph{embedding}
przed pozostałymi dwoma.
\end{note}

\medskip
Dwa dalsze hasła warto nazwać, bo są przypadkiem odwrotnym \dash{} miejscami,
w których wydanie polskie jest ustalone, a czytelnik może spodziewać się czegoś
innego. $\intcc{a}{b}$ zachowuje nawiasy kwadratowe zamiast ostrych, których uczy
polska szkoła, bo ISO 80000-2 i każdy artykuł, który czytelnik będzie otwierał,
piszą je właśnie tak. Żaden program się nad tym nie zatrzymuje, więc zapisem
tego jest Dodatek~\ref{app:B}, a nie pierwszy przedział, jaki czytelnik
napotyka. A tangens hiperboliczny zostaje $\tanh$ zamiast \emph{tgh}, co jest
testem klawiatury z sekcji~\ref{sec:D-notation} rozstrzygającym przeciw formie
polskiej, bo tę czytelnik będzie wpisywał.

\section{Dwie reguły, których nic nie sprawdzi}
\label{sec:D-untestable}
\index{terminology!translator rules}

Oba wydania są porównywane ze sobą linia po linii, a dwie reguły tłumacza leżą
poza wszystkim, co to porównanie widzi.

\textbf{Komentarze w listingu zostają angielskie.} Czytelnik jest szkolony do
czytania angielskiego kodu, a komentarz przetłumaczony na polski to jedyna
linia sesji, której nie dałoby się wkleić do terminala i wyszukać. Większość
listingów w tej książce to zapisy konsoli wygenerowane raz przez \code{code/}
i wciągane do obu wydań z tego samego pliku, więc ich komentarze są angielskie
z konstrukcji, a nie z czyjejś dyscypliny; listing pisany wprost w tekście
trzyma się tej reguły ręcznie.

\textbf{Zapisu stałej nigdy się nie tłumaczy.} $e$ to $e$, a $\pi$ to $\pi$
w obu wydaniach. To nazwy, a nie słowa, więc tłumaczenie którejkolwiek byłoby
tłumaczeniem identyfikatora \dash{} czyli znowu testem klawiatury z
sekcji~\ref{sec:D-notation}, zastosowanym do symbolu, który czytelnik wpisze,
a nie do funkcji, którą wywoła.
